%\documentclass[12pt]{scrreprt}
\documentclass[12pt]{report} 

% language may be romanian or english (default is english)
% type may be bachelor or master (default is bachelor)
\usepackage[language=romanian, type=bachelor]{style}

%\geometry{a4paper,top=2.5cm,left=3cm,right=2.5cm,bottom=2.5cm}
%in style
%controlling the appearance of your headers and footers
\usepackage{comment}
\usepackage{fancyhdr}
\setlength{\headheight}{15pt}
\pagestyle{fancy}
\lhead{}
\chead{}
\renewcommand{\headrulewidth}{0.2pt}
\renewcommand{\footrulewidth}{0.2pt}

\begin{document}

\specialization{Informatică}
\title{Sistem de Detectare și Răspuns la Intruziuni bazat pe eBPF}
\author{[Ovidiu-Mihai Moldovan]}
\supervisor{[Conf. univ. dr. Darius-Vasile Bufnea]}
				
\maketitle


\newpage
\thispagestyle{empty}
\mbox{}
\newpage
\pagenumbering{roman} 

\cleardoublepage
ABSTRACT
\vspace{0.5cm}	
\hrule
\vspace{0.5cm}	
%\cleardoublepage

This thesis presents the design and implementation of Sentinel, an eBPF-based Host Intrusion Detection and Response System (HIDS) for the Linux platform. Sentinel attaches kprobes to 20 security-relevant system calls, collects events through a 64 MiB ring buffer and processes them in a Go user-space daemon. A three-layer hybrid scoring engine assigns risk scores: static per-syscall weights calibrated on empirical malware data, a sliding-window burst detector, and an optional Isolation Forest ML scorer served through a REST API. Scores are mapped to six severity states (OK / WATCH / WARN / THROTTLE / ISOLATE / FREEZE) and a quarantine manager applies proportional containment actions via Linux cgroup v2 and nftables --- CPU throttling at 10\%, memory capping at 256 MiB, network isolation by cgroup inode, and full process freezing via SIGSTOP and \texttt{cgroup.freeze}. The system is validated against a reproducible suite of 20 attack scenarios covering kernel rootkits, privilege escalation, process injection, network exfiltration, persistence and fileless execution. Original contributions include the adaptive weighted scoring mechanism, the cgroup-inode-based selective network isolation, the PID-reuse-safe burst key, and the self-contained attack test suite.



\tableofcontents


\newpage
\pagenumbering{arabic}

\chapter{Introducere}
\label{intro}

\par Securitatea sistemelor de operare bazate pe Linux a dobândit o importanță tot mai mare pe măsură ce aceste sisteme au ajuns să susțină infrastructuri critice: servere web, platforme de containere, sisteme de calcul de înaltă performanță și dispozitive IoT. Creșterea numărului de atacuri direcționate împotriva nucleului Linux — rootkit-uri, tehnici de injecție de procese, escaladare de privilegii și execuție fără fișier pe disc — impune soluții de detecție capabile să opereze la nivelul cel mai profund al stivei software, fără a introduce latențe inacceptabile sau vulnerabilități suplimentare.

\par Lucrarea de față descrie proiectarea și implementarea sistemului \textit{Sentinel}, un sistem de detectare și răspuns la intruziuni bazat pe tehnologia \textit{extended Berkeley Packet Filter} (eBPF). Sistemul interceptează apeluri de sistem la nivel de kernel prin mecanismul kprobe, calculează un scor de risc compus pentru fiecare eveniment detectat și aplică automat acțiuni graduale de containment al proceselor suspectate. Sentinel funcționează fără a modifica codul kernel-ului și fără a fi necesară repornirea sistemului — proprietăți esențiale pentru medii de producție.


\section{Contextul și motivația}

\par Instrumentele tradiționale de monitorizare a securității, cum ar fi demonul de audit Linux \texttt{auditd} \cite{LinuxAudit} sau sistemul Snort bazat pe semnături de rețea \cite{Roesch1999}, au demonstrat limitări documentate în fața atacurilor sofisticate care operează exclusiv în memorie sau care manipulează structurile interne ale kernel-ului \cite{Cozzi2018}. Abordările bazate pe analiza jurnalelor suferă de o latență inerentă: evenimentul este produs, jurnalizat, scris pe disc și abia apoi analizat, permițând atacatorilor o fereastră de acțiune neobservată de ordinul secundelor sau minutelor.

\par eBPF oferă o paradigmă diferită. Programele eBPF se execută într-un sandbox verificat formal al kernel-ului Linux, atașate la puncte de instrumentare (kprobes, tracepoints, uprobes) cu o suprasarcină de ordinul zecilor de nanosecunde per eveniment \cite{Gregg2019}. Aceasta permite vizibilitate completă, în timp real, asupra comportamentului oricărui proces, fără modificarea kernel-ului sau instalarea de module de kernel suplimentare.

\par Complementar, metodele de detecție a anomaliilor bazate pe învățare automată — în special algoritmul Isolation Forest \cite{Liu2008IF} — s-au dovedit eficiente pentru identificarea comportamentelor atipice în spații de caracteristici de dimensionalitate medie, inclusiv în contextul securității informatice \cite{Buczak2016}. Integrarea unui scorer ML cu un agent de detecție bazat pe eBPF formează o arhitectură hibridă care combină viteza și determinismul regulilor statice cu capacitatea de generalizare a modelelor statistice.


\section{Enunțul problemei}

\par Problema adresată de această lucrare poate fi formulată astfel: \textit{cum poate fi proiectat un sistem de detectare și răspuns la intruziuni pentru Linux care (a) operează la nivel de kernel fără a-l modifica, (b) clasifică evenimentele de securitate în timp real pe baza unui scor de risc compus și (c) aplică automat acțiuni de containment proporționale cu severitatea detectată?}

\par Validarea empirică se realizează printr-o suită de 20 de scenarii de atac reproductibile, ce acoperă principalele categorii de amenințări documentate în literatura de specialitate: rootkit-uri de kernel, escaladare de privilegii, injecție de procese, exfiltrare de date în rețea, persistență și execuție de cod fără fișier pe disc.


\section{Obiectivele lucrării}

\par Obiectivele principale ale lucrării sunt:

\begin{itemize}
    \item Proiectarea și implementarea unui program eBPF pentru intercepția selectivă a apelurilor de sistem relevante din perspectiva securității, cu suprasarcină minimă asupra sistemului.
    \item Construirea unui motor de scoring hibrid care combină ponderi statice per syscall, detecție a activității în rafale (burst detection) și un scorer ML bazat pe Isolation Forest, care poate fi antrenat cu date actualizate fără înteruperi.
    \item Implementarea unui agent de carantinare care aplică acțiuni de izolare graduale — limitare CPU, izolare de rețea, înghețarea procesului — prin mecanismele cgroup v2 și nftables.
    \item Validarea întregului sistem printr-o suită de teste automatizate ce rulează 20 de scenarii de atac într-un mediu izolat bazat pe mașini virtuale Vagrant.
\end{itemize}


\section{Contribuții originale}

\par Contribuțiile originale ale lucrării sunt:

\begin{enumerate}
    \item \textbf{Un mecanism de scoring cu trei etape}: ponderi statice per apel de sistem (syscall), detectarea activității în rafale cu fereastră temporală și scorer Isolation Forest care poate fi antrenat cu date actualizate printr-un serviciu REST API, cu un mecanism de combinare ponderată a scorurilor care previne degradarea excesivă a scorului bazal.
    \item \textbf{O mașină de stări cu șase nivele} (OK / WATCH / WARN / THROTTLE / ISOLATE / FREEZE) care decuplează detecția de răspuns, permițând acțiuni proporționale cu severitatea detectată.
    \item \textbf{O suită de 20 de scenarii de atac reproductibile}, distribuite ca artefacte de testare, care acoperă vectori de atac documentați în literatura de specialitate — de la rootkit-uri tip Diamorphine până la tehnici de execuție de cod fără fișier și rootkit-uri eBPF de tip TripleCross.
    \item \textbf{Mecanismul de izolare de rețea bazat pe cgroup inode}: procesele carantinate sunt blocate în rețea prin reguli nftables care filtrează traficul după inodeul cgroup-ului, fără a afecta restul proceselor de pe sistem.
\end{enumerate}


\section{Structura lucrării}

\par Capitolul~\ref{chap:ch2} prezintă fundamentele teoretice: modelul de securitate Linux (apeluri de sistem, Linux Security Modules, namespaces și cgroups), arhitectura eBPF, clasificarea și desccrierea sistemelor de detectare a intruziunilor, bazele algoritmului Isolation Forest și o analiză comparativă a instrumentelor existente. Capitolul~\ref{chap:ch3} descrie analiza cerințelor și proiectarea detaliată a componentelor sistemului, inclusiv motorul de scoring și mașina de stări pentru carantinare. Capitolul~\ref{chap:ch4} acoperă implementarea concretă și rezultatele testării cu suita de 20 de atacuri. Capitolul~\ref{conclusions} sintetizează concluziile și direcțiile de cercetare viitoare.

\par \textit{Notă privind utilizarea instrumentelor de inteligență artificială}: la redactarea unor secțiuni ale acestei lucrări a fost utilizat asistentul Claude (Anthropic) pentru reformularea și îmbunătățirea lizibilității paragrafelor tehnice. De asemenea, asistentul Claude (Anthropic) a mai fost folosit în formatarea codului și în implementarea unor secțiuni de cod generice. Întregul conținut tehnic, implementarea software, excluzând secțiunile menționate anterior și rezultatele experimentale aparțin exclusiv autorului lucrării.

%\addcontentsline{toc}{chapter}{Introducere}
%\addcontentsline{toc}{chapter}{Introduction}

\chapter{Fundamente teoretice și contextul actual}
\label{chap:ch2}

\par Prezentul capitol stabilește fundamentul teoretic necesar înțelegerii și proiectării arhitecturii sistemului Sentinel. Sunt analizate mecanismele intrinseci de securitate ale nucleului Linux, inovațiile aduse de tehnologia eBPF, taxonomia sistemelor de detectare a intruziunilor, aparatul matematic al algoritmului Isolation Forest, precum și o analiză comparativă a soluțiilor consacrate în domeniu.


\section{Modelul de securitate Linux}
\label{sec:linux_security}

\subsection{Apeluri de sistem și frontiera kernel-utilizator}
\label{subsec:syscalls}

\par Granița fundamentală de izolare între codul cu privilegii reduse (spațiul utilizator / user-space) și codul cu privilegii absolute (nucleul / kernel-space) în sistemele Linux este definită de interfața apelurilor de sistem (syscalls). Orice acțiune care necesită interacțiunea cu resursele hardware sau logice ale sistemului --- instanțierea de conexiuni (\texttt{socket}), crearea de procese, manipularea sistemului de fișiere, alterarea privilegiilor --- trebuie să traverseze această interfață printr-o tranziție de context strict controlată \cite{Kerrisk2010}.

\par Din perspectiva securității cibernetice, monitorizarea apelurilor de sistem oferă o vizibilitate exhaustivă și imposibil de eludat asupra comportamentului real al unui proces: o rutină malițioasă care încearcă să încarce un modul de kernel va invoca inevitabil \texttt{finit\_module} (syscall 313), o tentativă de escaladare a privilegiilor va apela \texttt{setuid} (syscall 105) sau \texttt{capset} (syscall 126), în timp ce exfiltrarea datelor va fi tradusă prin generarea de apeluri \texttt{socket} și \texttt{connect} (syscall-urile 41, respectiv 42) \cite{Love2010}.

\par Kernel-ul Linux x86\_64 expune peste 300 de apeluri de sistem, dar numai un subansamblu specific prezintă interes direct din perspectiva securității. Sistemul Sentinel monitorizează 20 de apeluri de sistem, selectate pe baza incidenței lor empirice în anatomia malware-ului pentru Linux \cite{Cozzi2018}, acoperind categoriile: execuție de procese, modificare de privilegii, operații de memorie, activitate de rețea, operații pe fișiere sensibile și management de module de kernel. Tabelul~\ref{tab:syscalls_risk} prezintă o selecție cu ponderile de risc corespunzătoare.

\begin{table}[htbp]
\begin{center}
\begin{tabular}{|p{80pt}|p{40pt}|p{90pt}|p{60pt}|}
\hline
\textbf{Syscall} & \textbf{Nr.} & \textbf{Categorie alertă} & \textbf{Pondere} \\
\hline\hline
\texttt{finit\_module} & 313 & HiddenPID           & 95 \\
\hline
\texttt{init\_module}  & 175 & HiddenPID           & 95 \\
\hline
\texttt{delete\_module}& 176 & HiddenPID           & 85 \\
\hline
\texttt{ptrace}        & 101 & ProcessInjection    & 75 \\
\hline
\texttt{setuid}        & 105 & PrivilegeEscalation & 70 \\
\hline
\texttt{setgid}        & 106 & PrivilegeEscalation & 70 \\
\hline
\texttt{capset}        & 126 & PrivilegeEscalation & 70 \\
\hline
\texttt{execve}        & 59  & ParentChildAnomaly  & 50 \\
\hline
\texttt{mprotect}      & 10  & MemoryInjection     & 45 \\
\hline
\texttt{mmap}          & 9   & MemoryInjection     & 40 \\
\hline
\texttt{connect}       & 42  & NetworkAnomaly      & 15 \\
\hline
\texttt{socket}        & 41  & NetworkAnomaly      & 10 \\
\hline
\end{tabular}
\end{center}
\caption{Selecție de apeluri de sistem monitorizate de Sentinel cu ponderile de risc asociate}
\label{tab:syscalls_risk}
\end{table}


\subsection{Linux Security Modules}
\label{subsec:lsm}

\par Cadrul Linux Security Modules (LSM), teoretizat și implementat de Wright et al. \cite{Wright2002}, pune la dispoziție o arhitectură bazată pe cârlige de interceptare (hooks) plasate strategic în punctele decizionale ale nucleului. Această infrastructură permite modulelor de securitate (ex. SELinux, AppArmor) să impună politici stricte de tip Mandatory Access Control (MAC).

\par O limitare inerentă a modulelor LSM tradiționale este natura lor declarativă și statică, nefiind concepute pentru analiza comportamentală dinamică a secvențelor de execuție. Din acest motiv, un sistem de detectare a anomaliilor comportamentale, precum Sentinel, reprezintă o completare necesară și ortogonală față de mecanismele LSM. Deși eBPF permite atașarea programelor la punctele de interceptare LSM (prin \texttt{BPF\_PROG\_TYPE\_LSM}, funcționalitate introdusă în Linux 5.7), Sentinel utilizează tehnologia kprobes pentru a asigura o compatibilitate retroactivă mai largă pe nucleele $\geq 5.4$.


\subsection{Namespaces și grupuri de control (cgroups)}
\label{subsec:cgroups}

\par Mecanismele interne ale nucleului (PID, net, mount, user, UTS, IPC) asigură izolarea stării globale a kernel-ului în contexte virtualizate pentru grupuri distincte de procese, reprezentând fundamentul tehnologiilor moderne de containerizare \cite{Kerrisk2010}. În paralel, mecanismul Control Groups v2 (cgroup v2), fundamentat inițial de Menage \cite{Menage2007} și standardizat în nucleul 4.5, gestionează alocarea, limitarea și monitorizarea resurselor fizice (CPU, memorie, I/O bloc) pentru ierarhii de procese.

\par Spre deosebire de iterația anterioară (cgroup v1), arhitectura v2 propune o ierarhie unificată și un model de delegare mai coerent, eliminând ambiguitățile de configurare care apăreau la apartenența simultană a unui proces la sub-arbori de control diferiți. Arhitectura de izolare a sistemului Sentinel impune la pornire prezența controllerelor \texttt{cpu} și \texttt{memory} în \texttt{/sys/fs/cgroup/cgroup.controllers}, condiție satisfăcută nativ pe distribuțiile Debian 12 și Ubuntu 22.04+.


\section{Extended Berkeley Packet Filter (eBPF)}
\label{sec:ebpf}

\subsection{Evoluție de la BPF clasic la eBPF}
\label{subsec:ebpf_history}

\par Berkeley Packet Filter (BPF) a fost conceptualizat de McCanne și Jacobson în 1993, vizând dezvoltarea unei arhitecturi extrem de eficiente pentru capturarea și filtrarea traficului de rețea, operând printr-o mașină virtuală minimalistă pe 32 de biți \cite{McCanne1993}. Schimbarea de paradigmă pe care a propus-o BPF a fost executarea logicii de filtrare direct în kernel, evaluând pachetele la sursă pentru a preveni transferul inutil de date peste frontiera privilegiată.

\par Începând cu Linux 3.18 (2014), această arhitectură a suferit o extensie masivă sub acronimul eBPF: lărgirea registrelor la 64 de biți, îmbogățirea setului de instrucțiuni, introducerea structurilor de date partajate (maps) și capacitatea de atașare asincronă la aproape orice funcție a sistemului de operare \cite{Rice2023}. O inovație critică o reprezintă Verificatorul Formal (The Verifier), un mecanism care analizează codul înainte de execuție pentru a garanta finalizarea acestuia (lipsa buclelor infinite) și respectarea strictă a limitelor de memorie accesată. Compilatorul JIT (Just-In-Time) convertește ulterior bytecode-ul validat în cod mașină nativ, reducând la minimum costul de performanță (overhead).

\par Ecosistemul eBPF a evoluat rapid: bibliotecile \texttt{libbpf} și \texttt{cilium/ebpf} \cite{CiliumEBPF} oferă o interfață Go/C de nivel înalt pentru încărcarea și gestionarea programelor eBPF, iar CO-RE (Compile Once, Run Everywhere) rezolvă problema portabilității între versiuni de kernel prin relocări BTF.


\subsection{Maps eBPF și mecanismul ring buffer}
\label{subsec:ebpf_maps}

\par Maps eBPF sunt structuri de stocare persistente, expuse asincron atât spațiului kernel, cât și celui utilizator. Arhitectura Sentinel implementează trei structuri fundamentale:

\begin{itemize}
    \item \texttt{BPF\_MAP\_TYPE\_HASH}: tabelă de dispersie alocată pentru păstrarea excluderilor (\texttt{pid\_blacklist}).
    \item \texttt{BPF\_MAP\_TYPE\_PERCPU\_ARRAY}: vector instanțiat per-CPU pentru o contorizare atomică, fără blocaje (lock-free), a evenimentelor suprimate (\texttt{drop\_count}).
    \item \texttt{BPF\_MAP\_TYPE\_RINGBUF}: coadă circulară de mare viteză, dimensionată la 64 MiB, utilizată ca bus principal de transmitere a evenimentelor interceptate.
\end{itemize}

\par Ring buffer-ul (introdus în kernel 5.8) oferă avantaje față de perf buffer-ul precedent: elimină o copiere suplimentară de date, suportă rezervarea atomică a spațiului urmată de scriere directă și notificarea consumatorului printr-un descriptor de fișier abilitat în acest sens prin mecanismul epoll. Gregg \cite{Gregg2019} documentează că ring buffer-ul reduce latența de notificare cu 20--40\% față de perf buffer în scenariile de evenimente cu frecvență ridicată.


\subsection{Kprobes: instrumentarea dinamică a kernel-ului}
\label{subsec:kprobes}

\par Prin intermediul mecanismului kprobes (Kernel Probes), putem instrumenta codul nucleului în timp real, injectând rutine de tratare exact acolo unde avem nevoie, fără a atinge codul sursă și fără a recompila sistemul [CRKH05]. Funcționarea sa este elegantă: sonda înlocuiește discret instrucțiunea originală cu o întrerupere software (de tipul INT3 pe x86). Când execuția lovește această „capcană”, controlul este pasat instantaneu mașinii virtuale eBPF. După procesarea evenimentului, sistemul restaurează și execută instrucțiunea originală. Deși pare un proces complex, această deviere este extrem de rapidă, inducând o întârziere de doar 100–300 nanosecunde [VCP+20] – un cost insesizabil în contextul unui apel de sistem obișnuit.

\par Sentinelul atașează kprobe-uri pe funcțiile \texttt{\_\_x64\_sys\_\{name\}} corespunzătoare celor 20 de syscall-uri monitorizate. Atașarea este dinamică: la pornire, daemonul iterează peste toate programele din colecția eBPF cu prefix \texttt{kp\_} și le atașează la funcțiile kernel corespunzătoare, fără a fi necesară enumerarea lor explicită în codul Go.

\section{Sisteme de detectare a intruziunilor}
\label{sec:ids}

\subsection{Taxonomie și clasificare}
\label{subsec:ids_taxonomy}

\par Conceptul monitorizării securității a fost formalizat inițial de Anderson (1980) ca o metodă de prelucrare a jurnalelor de audit în scopul detectării utilizării neautorizate a sistemelor informatice \cite{Anderson1980}. Ulterior, Denning (1987) a publicat fundamentul teoretic al sistemelor IDS (sisteme de detecție a intruziunilor), propunând un model independent de sistem capabil să evalueze statistic deviațiile de la comportamentul istoric normal \cite{Denning1987}.

\par Taxonomia propusă de Axelsson (2000) \cite{Axelsson2000} oferă un cadru de clasificare în care Sentinel se raportează astfel:

\begin{itemize}
    \item \textit{Sursa datelor}: HIDS (Host-based) versus NIDS (Network-based). Sentinel este exclusiv un sistem HIDS, rezoluția sa fiind limitată la apelurile de sistem ale mașinii gazdă.
    \item \textit{Metoda de detecție}: Detecție bazată pe semnături versus anomalii. Sistemul de față abordează problema hibrid, combinând un set de euristici statice cu un model nesupervizat pentru anomalii.
    \item \textit{Frecvența analizei}: detecție continuă versus periodică. Sentinel operează în timp real.
    \item \textit{Tipul de răspuns}: alertare pasivă versus limitare activă (containment). Sentinel oferă un răspuns activ, autonom și gradual.
\end{itemize}


\subsection{Detecție bazată pe anomalii versus semnături}
\label{subsec:anomaly_vs_signature}

\par Abordarea bazată pe semnături excelează prin rata infimă a rezultatelor fals pozitive pentru atacurile cunoscute, însă suferă din cauza ineficienței totale în fața atacurilor neclasificate (zero-day) \cite{GarciaTe2009}. Alternativa, detecția anomaliilor, profilează starea de normalitate a sistemului; deși superioară în identificarea vectorilor de atac noi, aceasta este penalizată de un volum mai mare de rezultate fals pozitive \cite{Chandola2009}.

\begin{comment}
TODO: check factuality further
\par În stabilirea algoritmului adecvat, studiul lui Goldstein și Uchida \cite{Goldstein2016} demonstrează performanța net superioară a modelului Isolation Forest pe distribuții de date unde anomaliile reprezintă un procent extrem de redus ($< 5\%$), modelând perfect realitatea evenimentelor de securitate. Această alegere este validată și de lucrarea de sinteză elaborată de Buczak și Guven \cite{Buczak2016}, care evidențiază eficiența metodelor de ansamblu (ensemble methods) pe spații de date cu dimensionalitate variabilă.
\end{comment}

\section{Algoritmul Isolation Forest}
\label{sec:isolation_forest}

\par Isolation Forest (iForest), formulat de Liu, Ting și Zhou \cite{Liu2008IF, Liu2012IF}, schimbă paradigma clasică a detecției anomaliilor: în loc să profileze marea masă a datelor normale, algoritmul încearcă să izoleze explicit valorile aberante. Postulatul principal este că anomaliile reprezintă instanțe rare și divergente structural; prin urmare, o serie de partiționări aleatorii (random splits) va izola un punct anormal mult mai aproape de rădăcina unui arbore de decizie comparativ cu un punct benign.

\par Formal, dat fiind un set de antrenament $X = \{x_1, \ldots, x_n\} \subset \mathbb{R}^d$, construirea unui arbore de izolare presupune selecția aleatoare a unui atribut $q \in \{1, \ldots, d\}$ și partiționarea spațiului pe un punct $p$ ales dintr-o distribuție uniformă în intervalul $[\min_q(X), \max_q(X)]$. Evaluarea (scorul de anomalie) pentru un punct $x$ este obținută prin relația:

\begin{equation}
    s(x, n) = 2^{-\frac{E[h(x)]}{c(n)}}
    \label{eq:isolation_score}
\end{equation}

\noindent unde $E[h(x)]$ reprezintă valoarea așteptată (adâncimea medie) a lungimii căii către punctul $x$ generată pe ansamblul de arbori, iar $c(n) = 2H(n-1) - \frac{2(n-1)}{n}$ este adâncimea medie asimptotică într-un arbore binar de căutare nereușită, cu $H$ funcția armonică \cite{Liu2012IF}. Scorul variază în intervalul $(0, 1)$: valori aproape de $1$ indică anomalii puternice, valori aproape de $0.5$ indică puncte normale sau ambigue.

\par Deoarece scorul brut emis de model (valoarea funcției de decizie \texttt{decision\_function}, negativă pentru anomalii) necesită o conversie pe scala operațională $[0, 100]$, Sentinel aplică o transformare sigmoidă:

\begin{equation}
    \text{threat\_score} = \frac{100}{1 + e^{-25 \cdot (-\text{raw\_score})}}
    \label{eq:sigmoid_transform}
\end{equation}

\par Factorul de scalare 25 a fost ales empiric pentru a produce o separare clară între punctele normale (\texttt{threat\_score} $< 30$) și anomalii moderate (\texttt{threat\_score} $\in [50, 70]$), menținând în același timp sensibilitate pentru valorile extreme.

\par Complexitatea antrenării este $\mathcal{O}(t \cdot \psi \cdot \log \psi)$, cu $t$ numărul de arbori și $\psi$ dimensiunea sub-eșantionării (subsampling), în timp ce faza de inferență se execută în timp $\mathcal{O}(t \cdot \log \psi)$ per eveniment --- acceptabilă pentru bugetul de latență impus de constrângerile decizionale în timp real (RN2).


\section{Instrumente existente și lucrări corelate}
\label{sec:related_work}

\par \textbf{Falco} \cite{FalcoProject}, proiect incubat de CNCF și derivat din Sysdig, este cel mai apropiat echivalent open-source al Sentinelului: utilizează eBPF sau un driver de kernel pentru monitorizarea syscall-urilor și permite definirea de reguli de alertă în YAML. Spre deosebire de Sentinel, Falco nu implementează un scorer ML integrat, nu aplică acțiuni automate de containment (răspunsul este delegat unor sisteme externe precum Falcosidekick) și nu menține o bază de date locală a evenimentelor pentru analiza retrospectivă.

\par \textbf{auditd} \cite{LinuxAudit} este componenta de audit a kernel-ului Linux, standardizată și prezentă în toate distribuțiile majore. Jurnalizează apeluri de sistem pe baza unor reguli configurate static. Din pricina prelucrării asincrone și a interacțiunii constante cu sistemul de fișiere pentru stocarea jurnalelor, auditd induce o latență incompatibilă cu acțiunile de izolare în timp real, limitându-se la rolul de analiză retrospectivă (forensics).

\par \textbf{Virtual Machine Introspection (VMI)}: Metodologia teoretizată de Garfinkel și Rosenblum (2003) \cite{Garfinkel2003} propune securizarea prin mutarea monitorului în hipervizor, complet izolat de sistemul de operare monitorizat. Implementări precum DRAKVUF \cite{Lengyel2014} utilizează extensiile de virtualizare (Xen) pentru a analiza memoria mașinii virtuale din exterior. Limita severă a acestei abordări este dependența de virtualizarea hardware, fiind incompatibilă cu mediile bare-metal sau infrastructurile cloud moderne bazate pe containere.

\par \textbf{kGuard} \cite{Kemerlis2012} abordează o problemă complementară --- protecția kernel-ului împotriva atacurilor de tip return-to-user --- prin instrumentarea statică a codului kernel la compilare. Aceste abordări de securitate prin compilare sunt ortogonale față de Sentinel, care se concentrează pe detectarea comportamentului la runtime.

\par O constatare empirică importantă din literatura de specialitate provine din Cozzi et al. \cite{Cozzi2018}: 70\% din eșantioanele de malware Linux analizate utilizează socket-uri de rețea, 30\% apelează \texttt{ptrace} sau mapează pagini cu drepturi de execuție, iar 18\% utilizează module de kernel. Aceste statistici au constituit criteriul de selecție și de ponderare a syscall-urilor monitorizate de Sentinel.


\section{Peisajul amenințărilor în ecosistemul Linux și Cadrul MITRE ATT\&CK}
\label{sec:threat_landscape}

\par Evaluarea oricărui sistem modern de detectare a intruziunilor necesită raportarea riguroasă la un cadru standardizat de amenințări. În ecosistemul Linux, vectorii de atac au evoluat de la simple scripturi de exploatare la arhitecturi complexe, mapate exhaustiv în cadrul MITRE ATT\&CK \cite{MITRE_ATTCK}. Evaluarea eficacității sistemului Sentinel se fundamentează pe o analiză teoretică a 20 de scenarii de atac care acoperă tactici critice: persistență, escaladarea privilegiilor, evaziune defensivă, acces la credențiale și execuție.

\par În continuare sunt detaliate principalele categorii teoretice de amenințări care justifică arhitectura de monitorizare a sistemului:

\subsection{Rootkit-uri și persistența la nivel de nucleu}
\label{subsec:rootkits}

\par Compromiterea nucleului Linux s-a realizat istoric prin rootkit-uri de tip LKM (\textit{Loadable Kernel Module}), care vizează alterarea fluxului de execuție din spațiul privilegiat. Un etalon istoric este rootkit-ul \textit{Diamorphine}, apărut în 2013, care interceptează apeluri de sistem (precum \texttt{getdents64}) pentru a ascunde procese și asigură o portiță de acces (\textit{backdoor}) gestionată prin abuzarea apelului de sistem kill și interceptarea unor semnale POSIX predefinite pentru a declanșa funcții ascunse, care ar putea da privilegii atacatorului. Încărcarea acestor module este semnalizată primar de apelul \texttt{init\_module}.

\par Pentru a evita lăsarea de artefacte pe disc, atacatorii utilizează apelul \texttt{finit\_module}, care permite încărcarea unui modul compilat direct dintr-un descriptor de fișier menținut exclusiv în memoria volatilă. Amenințările de generație nouă includ rootkit-urile eBPF, exemplificate de modelul teoretizat de \textit{TripleCross} în 2022 \cite{TripleCross2022}. Acestea demonstrează capacitatea atacatorilor de a subverti însăși tehnologia eBPF pentru a intercepta apeluri de sistem și a orchestra injecții de cod invizibile scanărilor tradiționale.

\subsection{Escaladarea privilegiilor}
\label{subsec:privesc}

\par Obținerea drepturilor absolute de administrator (UID 0) pornind de la un cont restricționat reprezintă un obiectiv primar în faza de post-exploatare. Pe sistemele Linux, această acțiune este mediată în principal de apelurile \texttt{setuid} (pentru alterarea directă a identificatorului de utilizator) și \texttt{capset} (pentru modificarea fină a capabilităților POSIX). Un exemplu elocvent al severității acestei clase de atacuri este vulnerabilitatea \textit{PwnKit} (CVE-2021-4034), care a exploatat o corupție de memorie în componenta PolicyKit pentru a permite execuția cu succes a funcției \texttt{setuid(0)} pe majoritatea distribuțiilor majore \cite{PwnKit2022}.

\subsection{Evaziunea defensivă și execuția \textit{fileless}}
\label{subsec:fileless}

\par Malware-ul de tip \textit{fileless} execută instrucțiuni exclusiv în memorie, ocolind soluțiile antivirus tradiționale fundamentate pe analiza sistemului de fișiere. Această tehnică a fost facilitată structural de introducerea apelului \texttt{memfd\_create} în nucleul Linux 3.17 (2014), care permite crearea unui fișier anonim susținut doar de memoria RAM. Atacatorul poate scrie codul malițios (de exemplu, un binar ELF) în acest descriptor și îl poate executa prin apelul \texttt{execve} (referențiind calea virtuală \texttt{/proc/self/fd/}), fără ca executabilul să persiste vreodată pe un mediu de stocare fizic \cite{Gruss2016}.

\subsection{Injectarea de procese}
\label{subsec:process_injection}

\par Apelul de sistem \texttt{ptrace}, conceput inițial ca mecanism legitim pentru depanare (\textit{debugging}), a devenit instrumentul canonic pentru tehnicile de injectare a proceselor. Acesta permite unui proces malițios să suspende o țintă legitimă (\texttt{PTRACE\_ATTACH}), să îi altereze pointerul de instrucțiune (\texttt{PTRACE\_SETREGS}) și să scrie \textit{shellcode} direct în spațiul său de adresare (\texttt{PTRACE\_POKETEXT}).

\par O derivare complexă a acestei tehnici este golirea proceselor (\textit{Process Hollowing} sau \textit{RunPE}). Un proces legitim este instanțiat și suspendat, i se alocă o regiune de memorie nouă, simultan lizibilă, scriptibilă și executabilă (RWX) prin \texttt{mmap}, iar fluxul normal este înlocuit cu logica atacatorului \cite{Elastic2022}.

\subsection{Accesul la credențiale și controlul (command-and-control / c2)}
\label{subsec:c2}

\par \textbf{Furtul de credențiale.} Atacatorii enumeră sistematic fișiere sensibile (de exemplu, \texttt{/etc/shadow}, chei SSH private) printr-o cascadă de apeluri \texttt{openat}. Pentru extragerea parolelor aflate în uz, tehnici similare utilitarului Mimikatz (Windows) sunt aplicate pe Linux prin citirea memoriei proceselor rulante, mapată în \texttt{/proc/<pid>/mem} \cite{Bundles2021}.

\par \textbf{Exfiltrarea datelor.} Stabilirea canalelor de comandă și control implică adesea instanțierea de \textit{reverse shells}, procedeu în care gazda compromisă inițiază conexiuni TCP de ieșire (\textit{outbound}) către atacator, ocolind astfel regulile restrictive de firewall pentru traficul de intrare. Amprenta comportamentală constă în serii de apeluri \texttt{socket} și \texttt{connect}, urmate de deturnarea fluxurilor de intrare/ieșire prin invocarea unui shell via \texttt{execve} \cite{Cozzi2018}.
\chapter{Analiza și Proiectarea Sistemului Sentinel}
\label{chap:ch3}

\par Proiectarea sistemului Sentinel pornește de la necesitatea de a detecta și neutraliza amenințările cibernetice în timp real. Arhitectura propusă este modulară, decuplând logic instrumentarea nucleului (kernel), evaluarea scorului de risc și aplicarea măsurilor de izolare (containment). Cele trei componente majore interacționează printr-un flux de date unidirecțional: programul eBPF generează evenimente brute, daemonul Go le procesează și le persistă, iar managerul de carantinare le consumă pentru a aplica măsuri de restricție.


\section{Cerințe funcționale și nefuncționale}
\label{sec:requirements}

\par Cerințele funcționale principale care au ghidat proiectarea sunt:

\begin{itemize}
    \item \textbf{RF1}: Interceptarea la nivel de kernel a apelurilor de sistem relevante din perspectiva securității, fără a necesita recompilarea nucleului sau instalarea unor module terțe.
    \item \textbf{RF2}: Calcularea unui scor de risc compus pentru fiecare eveniment detectat, fundamentat pe cel puțin un set de reguli euristice statice.
    \item \textbf{RF3}: Clasificarea decizională a evenimentelor în șase stări de severitate, variind de la activitate benignă (OK) până la necesitatea blocării imediate a procesului (FREEZE).
    \item \textbf{RF4}: Aplicarea automată a acțiunilor de carantinare, strict proporționale cu severitatea detectată, utilizând mecanismele native și standardizate ale kernel-ului Linux.
    \item \textbf{RF5}: Persistarea evenimentelor și a acțiunilor de carantinare într-o bază de date locală, optimizată pentru audit și analiză retrospectivă.
    \item \textbf{RF6}: Expunerea unui serviciu REST API pentru integrarea opțională a unui modul de Machine Learning (ML), capabil de reantrenare dinamică fără întreruperea monitorizării.
\end{itemize}

\par Cerințele nefuncționale critice pentru un mediu de producție sunt:

\begin{itemize}
    \item \textbf{RN1}: Impactul operațional (overhead-ul) introdus de instrumentarea eBPF asupra procesorului (CPU) să fie limitat sub \textbf{5\%} într-un scenariu de activitate normală a sistemului.
    \item \textbf{RN2}: Latența de reacție --- măsurată de la momentul intercepției până la aplicarea efectivă a măsurii de carantinare (excluzând evaluarea ML) --- să nu depășească \textbf{100 ms}.
    \item \textbf{RN3}: Monitorizarea și raportarea continuă a eventualelor pierderi de evenimente cauzate de saturarea memoriei partajate (ring buffer).
\end{itemize}


\section{Arhitectura sistemului}
\label{sec:architecture}

\par Arhitectura sistemului Sentinel se bazează pe trei componente interdependente, ilustrate în schema fluxului de date de mai jos:

\begin{enumerate}
    \item \textbf{Programul eBPF} (\texttt{tracer.bpf.o}): rulează în spațiul privilegiat al nucleului, interceptează apelurile de sistem folosind sonde dinamice (\textit{kprobes}) și scrie evenimentele contextuale într-un \textit{ring buffer} eBPF de înaltă performanță.
    \item \textbf{Daemonul Go} (\texttt{sentinel-daemon}): operează în spațiul utilizator (user-space), consumă datele din ring buffer, evaluează riscul (prin motorul de scoring) și scrie alertele rezultate într-o bază de date SQLite (configurată în modul WAL).
    \item \textbf{Managerul de carantinare} (\texttt{sentinel-quarantine}): interoghează continuu baza de date pentru evenimente cu severitate critică și aplică direct acțiunile de izolare folosind \textit{cgroup v2} și \textit{nftables}.
\end{enumerate}

\begin{figure}[htbp]
\centering
\begin{verbatim}
  Kernel space                          User space
+------------------+   ring buffer   +-------------------+
|   eBPF kprobes   | --------------> |  sentinel-daemon  |
|  (tracer.bpf.o)  |                 | (scoring engine)  |
+------------------+                 +--------+----------+
                                              |  SQLite WAL
                                     +--------v----------+
[optional] <--- ML Scorer REST :8765 |sentinel-quarantine|
                                     | cgroup v2/nftables|
                                     +-------------------+
\end{verbatim}
\caption{Arhitectura sistemului Sentinel: fluxul de date de la kprobe la carantinare.}
\label{fig:architecture}
\end{figure}

\begin{comment}
TODO: check rellevance
\par Decuplarea daemonului de managerul de carantinare, realizată prin intermediul bazei de date SQLite, oferă un avantaj major de reziliență: permite reluarea procesării evenimentelor din punctul ultimei opriri în cazul unui crash (toleranță la defecte) și elimină dependența de sincronizare strictă între producerea evenimentului și aplicarea pedepsei. O latență de interogare asincronă de 500 ms a bazei de date este considerată un compromis acceptabil pentru majoritatea acțiunilor de izolare, menținând în același timp un impact minim asupra I/O-ului sistemului.
\end{comment}

\section{Proiectarea programului eBPF}
\label{sec:ebpf_design}

\par Componenta eBPF este proiectată să asigure interceptarea a 20 de apeluri de sistem relevante. Fiecare apel este instrumentat printr-un punct de interceptare kprobe dedicată, având prefixul \texttt{kp\_}. Arhitectura folosește două tipologii principale de rutine de tratare (handlers):

\begin{itemize}
    \item \texttt{submit\_event(nr)}: funcție generică utilizată pentru apelurile care nu operează cu căi de fișiere. Extrage un set esențial de metadate: marca temporală (nanosecunde), PID, UID, numărul apelului și denumirea executabilului (\texttt{comm}).
    \item \texttt{submit\_event\_with\_path(nr, ctx)}: variantă extinsă, utilizată pentru apelurile care manipulează fișiere (ex. \texttt{openat}, \texttt{unlinkat}, \texttt{renameat2}). Această rutină utilizează asistentul \texttt{bpf\_probe\_read\_user\_str} pentru a extrage dinamic pointer-ul către șirul de caractere (din registrul \texttt{rsi} al structurii \texttt{pt\_regs}) peste granița spațiului utilizator.
\end{itemize}

\par Handler-ul pentru \texttt{write} (syscall 1) necesită un tratament special. Deoarece argumentul său este un descriptor de fișier (număr întreg), nu o cale, handler-ul eBPF ignoră explicit scrierile către descriptorii standard (0, 1 și 2) --- acestea generează un volum ridicat de evenimente cu relevanță de securitate scăzută. Pentru descriptorii rămași, funcția codifică valoarea numerică într-un șir de forma \texttt{"fd:N"}. Sarcina traducerii este delegată daemonului Go din spațiul utilizator, care interoghează pseudo-sistemul \texttt{/proc/\{pid\}/fd/\{N\}} pentru a obține calea absolută, menținând astfel execuția eBPF fluidă \cite{Kerrisk2010}.

\par Filtrul pe tabela de dispersie \texttt{pid\_blacklist} previne fenomenul de auto-instrumentare (detectia propriilor apeluri de sistem): imediat după lansare, daemonul Go își scrie propriul PID în această structură partajată de tip dicționar. La fiecare kprobe, codul eBPF verifică această listă și abandonează silențios interceptările generate de componentele de securitate proprii, eliminând astfel bucla de feedback ce ar satura ring buffer-ul.


\section{Proiectarea motorului de scoring}
\label{sec:scoring_design}

\par Motorul de scoring rafinează progresiv nivelul de amenințare printr-un sistem decizional etapizat, generând un scor final în intervalul $[0, 100]$.

\subsection{Etapa 1: Evaluarea statică a apelurilor de sistem}
\label{subsec:static_weights}

\par Fiecărui apel interceptat îi este asociată o pondere euristică $w_{nr} \in [5, 95]$, calibrată empiric conform studiilor privind vectorii de atac Linux \cite{Cozzi2018, Buczak2016}. Ponderile extreme (95) sunt alocate acțiunilor care nu au utilitate legitimă normală (ex. \texttt{init\_module}), în timp ce valorile minime (5) corespund apelurilor zgomotoase (\texttt{openat}). Dacă procesul inițiator beneficiază de privilegii administrative ($uid = 0$), ponderea este majorată cu un factor de 1,3, reflectând gravitatea exponențială a atacurilor executate cu drepturi depline:

\begin{equation}
    \text{score}_{\text{static}}(nr, uid) = w_{nr} \cdot \begin{cases} 1.3 & \text{dacă } uid = 0 \\ 1.0 & \text{altfel} \end{cases}
    \label{eq:static_score}
\end{equation}

\par Amplificarea pentru procesele root reflectă faptul că aceleași syscall-uri au un impact semnificativ mai mare când sunt executate cu privilegii complete.


\subsection{Etapa 2: Analiza comportamentală a activității în rafală (Burst Detection)}
\label{subsec:burst_detection}

\par Atacurile de tip forță brută (SSH brute-force) sau epuizare a resurselor (fork bombs) se manifestă prin frecvențe anormale de apeluri identice. Etapa 2 implementează o fereastră temporală glisantă (sliding window) care monitorizează agregarea pe tripletul (\texttt{PID}, \texttt{syscall\_nr}, \texttt{start\_time}). Includerea timpului de inițializare a procesului (\texttt{start\_time}) previne coliziunile cauzate de fenomenul de \textit{PID wrap-around} în sistemele cu rotație rapidă a proceselor.

\begin{table}[htbp]
\begin{center}
\begin{tabular}{|p{90pt}|p{50pt}|p{70pt}|p{55pt}|p{60pt}|}
\hline
\textbf{Syscall(uri)} & \textbf{Prag} & \textbf{Fereastră} & \textbf{Cooldown} & \textbf{Scor burst} \\
\hline\hline
\texttt{socket}, \texttt{connect} & 5  & 3 s  & 10 s & 75.0 \\
\hline
\texttt{clone}                     & 20 & 2 s  & 10 s & 35.0 \\
\hline
\texttt{openat}                    & 8  & 3 s  & 10 s & 60.0 \\
\hline
\end{tabular}
\end{center}
\caption{Parametrii de detectare a activității în rafale}
\label{tab:burst_params}
\end{table}

\par Intervalul de cooldown de 10 s previne fenomenul de \textit{alert fatigue} --- odată emis un alert pentru un triplet (PID, syscall), acesta nu va mai fi reeditat timp de 10 secunde pentru aceeași manifestare continuă, reducând zgomotul în baza de date.


\subsection{Etapa 3: Analiza contextuală bazată pe căile fișierelor}
\label{subsec:path_alerts}

\par A treia etapaă realizează o potrivire de tipare (pattern matching) pe căile fișierelor accesate de syscall-urile relevante (\texttt{write} cu fd rezolvat, \texttt{openat}, \texttt{unlinkat}, \texttt{renameat2}), vizând tactici cunoscute de post-exploatare:

\begin{itemize}
    \item \textbf{Log Tamper} (scor 50): scriere sau alterare a directoarelor \texttt{/var/log/}, specifică tentativelor de ștergere a urmelor.
    \item \textbf{LSMFileWrite} (scor 50): manipularea fișierelor \texttt{/etc/cron*} (asigurarea persistenței) sau deschiderea bibliotecilor dinamice (\texttt{.so}) din directoare mapate în memorie (\texttt{/tmp/}, \texttt{/dev/shm/}), un indicator pentru executarea \textit{fileless malware}.
    \item \textbf{Privilege Escalation} (scor 65): alterarea directă a bazelor de date de autentificare (\texttt{/etc/passwd}, \texttt{/etc/shadow}, \texttt{/etc/sudoers}).
\end{itemize}

\par Dacă o regulă contextuală este validată, scorul acesteia \textbf{suprascrie} scorul static inițial (în loc să se cumuleze prin adunare), prevenind astfel inflația artificială a riscului și reflectând cu acuratețe severitatea tacticii interceptate.


\subsection{Integrarea scorerului ML}
\label{subsec:ml_integration}

\par Dacă URL-ul scorerului ML este configurat la pornire, daemonul Go apelează asincron endpoint-ul \texttt{POST /score} pentru fiecare eveniment cu scor $\geq 10$. Apelul este condiționat de o limită (\textit{timeout}) de 50 ms; în absența unui răspuns în timp util, evaluarea continuă exclusiv static. Scorul final ia forma unei fuziuni ponderate (Score Blending):

\begin{equation}
    \text{score}_{\text{final}} = (1 - w) \cdot \text{score}_{\text{static}} + w \cdot \text{score}_{\text{ML}}
    \label{eq:blended_score}
\end{equation}

\noindent unde $w \in \{0.1, 0.2, 0.5, 0.7\}$ variază invers proporțional cu gravitatea scorului static:

\begin{itemize}
    \item $w_{nr} \geq 85$: $w = 0.1$ --- pentru vectori destructivi clari, logica statică domină absolut decizia (90\%).
    \item $w_{nr} \in [60, 85)$: $w = 0.2$ --- analiza statică păstrează o autoritate de 80\%.
    \item $w_{nr} \in [30, 60)$: $w = 0.5$ --- decizie echilibrată (50/50) între reguli și comportamentul algoritmic.
    \item $w_{nr} < 30$: $w = 0.7$ --- sistemul ML primește libertate majoritară de a discerne anomaliile fine în zgomotul apelurilor comune.
\end{itemize}

\par Un floor de $0.5 \cdot \text{score}_{\text{static}}$ garantează că modulul ML nu poate submina evaluarea de securitate la mai puțin de jumătate din valoarea sa originară, prevenind atacatorii să camufleze exploit-uri severe în tipare artificiale de \textit{trafic normal}.


\section{Mașina de stări pentru carantinare}
\label{sec:state_machine}

\par Severitatea finală este tradusă operațional într-una din cele șase stări acționabile, ordonate crescător: OK $\rightarrow$ WATCH $\rightarrow$ WARN $\rightarrow$ THROTTLE $\rightarrow$ ISOLATE $\rightarrow$ FREEZE, conform pragurilor din Tabelul~\ref{tab:states}.

\begin{table}[htbp]
\begin{center}
\begin{tabular}{|p{65pt}|p{70pt}|p{175pt}|}
\hline
\textbf{Stare} & \textbf{Prag scor} & \textbf{Acțiune aplicată} \\
\hline\hline
OK       & $< 30$    & Eveniment considerat benign și filtrat; nicio acțiune \\
\hline
WATCH    & $\geq 30$ & Eveniment suspect, jurnalizat pentru auditare în baza de date \\
\hline
WARN     & $\geq 50$ & Alertare activă; eveniment propagat în dashboard \\
\hline
THROTTLE & $\geq 70$ & Limitare CPU la 10\% prin directiva \texttt{cpu.max} în cgroup v2 \\
\hline
ISOLATE  & $\geq 85$ & Limitare de memorie la 256 MiB și izolare completă a rețelei via nftables \\
\hline
FREEZE   & $\geq 95$ & Suspendarea procesului la nivel de kernel (\texttt{SIGSTOP} și \texttt{cgroup.freeze}); autorecuperare programată după 5 minute \\
\hline
\end{tabular}
\end{center}
\caption{Mașina de stări: praguri de scor și acțiuni de carantinare corespunzătoare}
\label{tab:states}
\end{table}

\par Aceste praguri au fost aliniate operațional: un apel banal executat de contul de sistem (\texttt{openat} ca root, scor $6.5$) este ignorat; o rafală nejustificată de conexiuni atinge nivelul THROTTLE; în timp ce simpla încercare de instalare a unui modul extern (\texttt{finit\_module} ca root, scor $123.5$, plafonat la $95$) atrage suspendarea absolută (FREEZE) a procesului atacator.

\par Izolarea rețelei în starea ISOLATE se realizează prin nftables: managerul de carantinare creează un tabel dedicat \texttt{inet sentinel\_quarantine} și aplică reguli restrictive de filtrare (drop) bazate pe identificatorul (inode-ul) atribuit de cgroup procesului compromis \cite{Menage2007}. Această granulare asigură restricționarea chirurgicală a ramurii malițioase, fără a afecta comunicațiile serverului de producție.


\section{Schema bazei de date}
\label{sec:db_schema}

\par Sentinelul utilizează SQLite în modul WAL (Write-Ahead Logging) cu parametrul activ de gestionare a blocajelor (\texttt{\_busy\_timeout=5000}). Această configurație tranzacțională suportă citiri recurente (efectuate de managerul de carantinare) care nu invalidează și nu blochează fluxul continuu de scriere asigurat de daemonul principal. SQLite WAL permite citiri concurente cu o scriere activă, satisfăcând astfel modelul de acces al sistemului.

\par Sunt menținute două tabele principale:

\begin{itemize}
    \item \texttt{events}: jurnalizează totalitatea evenimentelor relevante. Câmpurile includ \texttt{alert\_type}, \texttt{score}, \texttt{state}, \texttt{pid}, \texttt{comm}, \texttt{timestamp} (ISO 8601) și un payload flexibil de metadate \texttt{detail} (în format JSON, conținând câmpurile \texttt{ts\_ns}, \texttt{syscall\_name}, \texttt{fname} și \texttt{source}).
    \item \texttt{quarantine\_log}: o pistă de audit irefutabilă pentru măsurile de remediere automată. Leagă acțiunea de izolare de declanșatorul inițial prin cheia externă \texttt{event\_id} și înregistrează metadatele restricției aplicate: acțiunile sistemice (\texttt{actions} codificat JSON), confirmarea succesului, erorile OS și durata temporală de execuție a pedepsei (\texttt{latency\_ms}).
\end{itemize}
 
\par Câmpul \texttt{latency\_ms} din \texttt{quarantine\_log} înregistrează durata totală de aplicare a acțiunilor de containment, de la selectarea evenimentului până la confirmarea acțiunilor, oferind un mecanism integrat prin care performanța de reacție a sistemului poate fi evaluată și validată obiectiv în concordanță cu cerința nefuncțională RN2.
\chapter{Implementare și Testare}
\label{chap:ch4}

\par Capitolul de față detaliază implementarea arhitecturii Sentinel și evaluează eficacitatea sistemului printr-o suită exhaustivă de 21 de scenarii de atac. Fiecare componentă este analizată individual, cu accent pe deciziile arhitecturale netriviale, optimizările de performanță și aspectele de siguranță a execuției care fundamentează corectitudinea sistemului.


\section{Mediul de dezvoltare și testare}
\label{sec:dev_env}

\par Dezvoltarea a fost realizată pe un host Ubuntu 22.04 (kernel 6.17.0). Testarea se desfășoară într-o mașină virtuală Vagrant cu Debian 12 (kernel 6.1.0-amd64, cgroup v2 activat implicit), izolând atacurile de mediul de dezvoltare. Mediul virtualizat este o necesitate absolută: scenariile REP-17 și REP-18 implică injectarea de module de kernel corupte, iar scenariul REP-01 utilizează un rootkit real de tip Diamorphine --- rularea acestora pe sistemul gazdă ar duce inevitabil la compromiterea totală a acestuia sau la instabilitate severă (Kernel Panic).

\par Stiva de unelte utilizată:

\begin{itemize}
    \item \textbf{clang 14+} cu ținta \texttt{bpf} pentru compilarea programului eBPF; flag-ul \texttt{-D\_\_TARGET\_ARCH\_x86} selectează macro-urile corecte de acces la registrele procesorului.
    \item \textbf{Go 1.21+} cu biblioteca \texttt{github.com/cilium/ebpf} \cite{CiliumEBPF} pentru implementarea daemonului de monitorizare și a managerului de carantinare.
    \item \textbf{Python 3.12} cu scikit-learn \cite{Pedregosa2011}, FastAPI și uvicorn pentru motorul de inferență ML și expunerea serviciului REST.
    \item \textbf{SQLite 3.40+} în modul WAL (Write-Ahead Logging) pentru persistența concurentă a datelor.
    \item \textbf{Vagrant 2.4+} cu provider VirtualBox pentru orchestrarea mediului de test.
\end{itemize}


\section{Implementarea programului eBPF}
\label{sec:ebpf_impl}

\par Programul eBPF (\texttt{tracer.bpf.c}, 157 de linii) este compilat cu comanda:

\begin{verbatim}
clang -target bpf -O2 -g -D__TARGET_ARCH_x86 \
  -I/usr/include -I/usr/include/x86_64-linux-gnu \
  -c tracer.bpf.c -o tracer.bpf.o
\end{verbatim}

\par Flag-ul \texttt{-O2} este esențial din două motive: verificatorul eBPF impune o limită de 1 milion de instrucțiuni per execuție și respinge programele cu structuri de control prea complexe; în plus, backend-ul LLVM pentru arhitectura BPF necesită un arbore de sintaxă optimizat pentru a emite \textit{opcodes} BPF valide --- codul neoptimizat (\texttt{-O0}) generează frecvent accesări de memorie pe care Verificatorul le respinge ca nesigure.

\par O decizie de implementare importantă privește extragerea sigură a căilor absolute ale fișierelor. Handler-ul \texttt{submit\_event\_with\_path} utilizează secvența:

\begin{verbatim}
struct pt_regs *inner = (struct pt_regs *)PT_REGS_PARM1(ctx);
unsigned long fname_uptr = 0;
bpf_probe_read_kernel(&fname_uptr, sizeof(fname_uptr), &inner->rsi);
bpf_probe_read_user_str(e->fname, sizeof(e->fname),
                        (const char *)fname_uptr);
\end{verbatim}

\par Kprobe-urile sunt atașate la wrapper-urile \texttt{\_\_x64\_sys\_*} care primesc un singur argument --- un pointer la structura \texttt{pt\_regs} cu argumentele originale ale syscall-ului. Accesul direct la argumentele syscall-ului necesită dereferențierea unui nivel suplimentar de \texttt{pt\_regs} (\textit{inner regs}), o idiomă documentată în \cite{Rice2023} și necesară pentru compatibilitatea cu versiunile kernel $\geq 4.17$ pe x86\_64.

\par Mecanismul de drop count este implementat cu \texttt{\_\_sync\_fetch\_and\_add}, operație atomică per-CPU: când ring buffer-ul este plin și rezervarea eșuează, contorul de drop al CPU-ului curent este incrementat. Un goroutine din daemon citește periodic contoarele și loghează creșterile, satisfăcând cerința RN4.


\section{Implementarea daemonului Go}
\label{sec:daemon_impl}

\par Daemonul Go (\texttt{daemon/main.go}, 497 de linii) orchestrează pipeline-ul de detecție. La pornire, daemonul:

\begin{enumerate}
    \item Ridică limita \texttt{RLIMIT\_MEMLOCK} prin \texttt{rlimit.RemoveMemlock()}, necesară pentru alocarea ring buffer-ului de 64 MiB.
    \item Blochează thread-ul OS cu \texttt{runtime.LockOSThread()} --- această operație este mandatorie, deoarece anumite apeluri de sistem eBPF și asocieri de namespace-uri sunt dependente de contextul specific al thread-ului care le inițiază \cite{CiliumEBPF}.
    \item Încarcă colecția eBPF din fișierul \texttt{.bpf.o}, atașează kprobe-urile și deschide reader-ul de ring buffer.
    \item Inserează propriul PID în \texttt{pid\_blacklist} pentru a preveni buclele de auto-instrumentare.
\end{enumerate}

\par Bucla principală citește din ring buffer în mod blocant (\texttt{rb.Read()}), parsează evenimentele prin \texttt{binary.Read} (little-endian, pentru a respecta ordinea octeților x86\_64) și aplică secvențial cele trei straturi de scoring. Operațiunea de inserare în SQLite utilizează o instrucțiune pregătită (\textit{prepared statement}) pentru a elimina overhead-ul de parsare SQL per eveniment.

\par Detectarea rafale (\texttt{burstTracker}) menține o mapă \texttt{windows[key][]int64} cu timestamp-urile apelurilor dintr-o fereastră glisantă. Cheia de burst este:

\begin{equation}
    \text{key}(pid, nr) = (\text{pid} \ll 32) \mid (nr \ll 16) \mid (\text{start\_time} \mathrel{\&} 0\text{xFFFF})
    \label{eq:burst_key}
\end{equation}

\par Incluzând \texttt{start\_time} (câmpul 22 din \texttt{/proc/\{pid\}/stat}, exprimat în \textit{ticks} de ceas) în cheie, sistemul distinge între procese diferite care rulează cu același PID pe parcursul testelor --- o situație frecventă în suita de atacuri, unde procesele au durată scurtă de viață. Utilizarea exclusivă a ultimilor 16 biți din \texttt{start\_time} introduce o probabilitate de coliziune matematic neglijabilă ($\approx 1/65536$), un compromis excelent între precizia detecției și amprenta de memorie.

\par Un goroutine separat curăță intrările expirate la fiecare 30 de secunde pentru a preveni creșterea nelimitată a memoriei. Altul monitorizează contorul de drop la fiecare 10 secunde.


\section{Implementarea scorerului ML}
\label{sec:scorer_impl}

\par Scorerul ML (\texttt{scorer/isolation\_forest.py}, 195 de linii) expune un API FastAPI cu trei endpoint-uri: \texttt{GET /health}, \texttt{POST /score} și \texttt{POST /model/train}. Modelul este încărcat \textit{lazy} la primul apel \texttt{/score}, eliminând latența de inițializare la pornirea serviciului.

\par Faza de inginerie a caracteristicilor (\textit{feature engineering}) transformă telemetria brută într-un vector de 26 de trăsături. Primele 21 constituie un encoding one-hot al syscall-ului (câte un bit per syscall monitorizat, conform indexului \texttt{SYSCALL\_INDEX}). Caracteristicile suplimentare sunt:

\begin{itemize}
    \item $v_{21}$: flag binar $\mathbb{1}[uid = 0]$ --- capturează comportamentul proceselor root.
    \item $v_{22}$: ponderea de risc normalizată $w_{nr}/100$ --- injectează cunoașterea a priori a riscului syscall-ului în algoritmul nesupervizat. Isolation Forest nu poate distinge semantic între un apel inofensiv și unul critic; adăugarea acestei greutăți forțează algoritmul să izoleze spațial evenimentele inerent periculoase.
    \item $v_{23}$: ora din zi normalizată $h/23$ --- capturează activitate nocturnă atipică.
    \item $v_{24}$: flag weekend $\mathbb{1}[\text{weekday} \geq 5]$.
    \item $v_{25}$: $\log_{10}(\text{pid})/6$ --- diferențiază procesele de sistem (PID mic) de procesele spațiului utilizator.
\end{itemize}

\par Modelul Isolation Forest este antrenat pe evenimentele cu scor static $< 30$ din baza de date (care reprezintă comportamentul de bază), cu parametrii $n\_estimators=200$ și $contamination=0.01$. Vectorii sunt standardizați prin \texttt{StandardScaler} \cite{Pedregosa2011} înaintea antrenării, iar scalarea este aplicată identic la inferență, utilizând parametrii din faza de antrenament. Modelul antrenat este persistat prin \texttt{joblib} și poate fi reantrenat cu date actualizate prin \texttt{POST /model/train} fără a opri serviciul.

\par Conversia scorului brut IF (valoarea \texttt{decision\_function}, negativă pentru anomalii, pozitivă pentru puncte normale) în scor de amenințare pe scala 0--100 utilizează transformarea sigmoidă cu coeficientul de creștere 25, care asigură o tranziție rapidă (abruptă) între stările benigne și maligne --- formula~\ref{eq:sigmoid_transform} din Capitolul~\ref{chap:ch2}. Un punct cu \texttt{raw\_score} $= -0.1$ (ușor anormal) obține \texttt{threat\_score} $\approx 92$, în timp ce un punct cu \texttt{raw\_score} $= 0.1$ (ușor normal) obține \texttt{threat\_score} $\approx 8$.


\section{Implementarea managerului de carantinare}
\label{sec:quarantine_impl}

\par Managerul de carantinare (\texttt{quarantine/main.go}, 306 de linii) validează la pornire disponibilitatea cgroup v2 prin verificarea magic number-ului sistemului de fișiere \texttt{/sys/fs/cgroup} (\texttt{0x63677270} = ``cgrp'' în ASCII) și prezența controllerelor \texttt{cpu} și \texttt{memory} în \texttt{/sys/fs/cgroup/cgroup.controllers}. Această verificare este esențială: pe sisteme cu cgroup v1 sau cgroup v2 fără controllerele activate, acțiunile de carantinare ar eșua silențios.

\par Acțiunile de izolare sunt implementate pe trei niveluri de asprime:

\begin{itemize}
    \item \textbf{THROTTLE}: Creează cgroup \texttt{/sys/fs/cgroup/sentinel/quarantine\_\{pid\}} prin \texttt{os.MkdirAll}, scrie PID-ul procesului în \texttt{cgroup.procs} și scrie \texttt{"10000 100000"} în \texttt{cpu.max}. Această configurare permite procesului să utilizeze 10000 microsecunde din fiecare perioadă de 100000 microsecunde, echivalent cu 10\% CPU --- suficient pentru a sugruma operațiunile intensive (ex. \textit{cryptojacking}), fără a opri complet procesul.
    \item \textbf{ISOLATE}: Creează cgroup, scrie limita de memorie de $256 \times 2^{20}$ octeți în \texttt{memory.max}, apoi aplică regulile nftables. Valoarea de 256 MiB a fost aleasă ca suficientă pentru a permite terminarea normală a procesului, dar insuficientă pentru exfiltrarea unor volume mari de date în memorie. Izolarea de rețea utilizează inodul cgroup-ului ca identificator de filtrare în nftables, evitând necesitatea de network namespaces (care nu pot fi aplicate unui proces deja în execuție fără cooperarea sa).
    \item \textbf{FREEZE}: Trimite \texttt{SIGSTOP} procesului (oprire imediată, indiferent de starea actuală), creează cgroup-ul și scrie \texttt{"1"} în \texttt{cgroup.freeze}. Combinația SIGSTOP + \texttt{cgroup.freeze} este critică: chiar dacă un atacator extern ar încerca să reia procesul printr-un semnal \texttt{SIGCONT}, limitarea impusă la nivel de kernel de cgroup previne planificarea (\textit{scheduling}) procesului de către procesor.
\end{itemize}

\par Izolarea de rețea pentru starea ISOLATE utilizează inodul cgroup-ului ca identificator de filtrare în nftables. Inodul se obține prin \texttt{syscall.Stat} pe calea cgroup-ului; nftables suportă filtrarea \texttt{meta cgroup} prin inod în kernel $\geq 5.10$. Această abordare este net superioară izolării prin \textit{network namespaces}, deoarece poate fi aplicată ``la cald'', fără a necesita repornirea sau cooperarea procesului malițios.

\par Auto-dezghețarea pentru starea FREEZE este gestionată de un goroutine \texttt{unfreezeWatcher} care verifică la fiecare 5 secunde procesele cu deadline expirat. Dezghețarea constă în scrierea \texttt{"0"} în \texttt{cgroup.freeze}, trimiterea \texttt{SIGCONT} și eliminarea procesului din cgroup. Intervalul implicit de 5 minute este configurabil la pornire prin flag-ul \texttt{--unfreeze-after}.


\section{Suita de teste de atac}
\label{sec:attack_suite}

\subsection{Metodologia de testare}
\label{subsec:test_method}

\par Fiecare scenariu de atac este definit printr-un fișier \texttt{attack.toml} cu patru câmpuri: \texttt{id}, \texttt{name}, \texttt{alert\_type} (tipul de alertă așteptat) și \texttt{expected\_state} (starea minimă așteptată). Scriptul \texttt{run.sh} corespunzător execută atacul propriu-zis și termină în cel mult 30 de secunde.

\par Orchestratorul Python \texttt{run\_sentinel.py} automatizează ciclul de testare, garantând izolarea mediului între teste prin golirea jurnalurilor WAL și restaurarea bazei de date la starea de zero:

\begin{enumerate}
    \item Resetare baza de date (\texttt{DELETE FROM events; PRAGMA wal\_checkpoint}).
    \item Execuție \texttt{run.sh} în mașina virtuală prin \texttt{vagrant ssh}.
    \item Așteptare flux de evenimente ($t_{sleep}$ configurabil per scenariu).
    \item Extragere snapshot baza de date din VM prin SCP.
    \item Verificare: există cel puțin un eveniment kprobe (nu sintetic) cu tipul și scorul corecte?
\end{enumerate}

\par Un test trece dacă și numai dacă există cel puțin un eveniment cu sursă \texttt{"kprobe"} (câmpul \texttt{detail.source}) cu \texttt{alert\_type} corespunzător și \texttt{score} $\geq$ pragul minim. Cerința de evenimente reale (kprobe) distinge testele cu detecție reală de eventualele injecții sintetice de testare.

\subsection{Scenariile de bază (REP-01 -- REP-12)}
\label{subsec:basic_attacks}

\par Setul de bază cuprinde 11 scenarii, prezentate în Tabelul~\ref{tab:basic_attacks}, alese să acopere categoriile principale de amenințări identificate de Cozzi et al. \cite{Cozzi2018}.

\begin{table}[htbp]
\begin{center}
\begin{tabular}{|p{45pt}|p{110pt}|p{75pt}|p{55pt}|}
\hline
\textbf{ID} & \textbf{Scenariu} & \textbf{Tip alertă} & \textbf{Stare} \\
\hline\hline
REP-01 & Rootkit Diamorphine                   & HiddenPID           & FREEZE   \\
\hline
REP-02 & Escaladare privilegii (pkexec/setuid) & PrivilegeEscalation & WARN     \\
\hline
REP-04 & Exfiltrare date (netcat)              & NetworkAnomaly      & THROTTLE \\
\hline
REP-05 & Webshell (execve din web server)      & ParentChildAnomaly  & WARN     \\
\hline
REP-06 & Persistență prin cron                & LSMFileWrite        & WARN     \\
\hline
REP-07 & Mapare memorie RWX (mmap)             & MemoryInjection     & WARN     \\
\hline
REP-08 & Manipulare fișiere log               & LogTamper           & WARN     \\
\hline
REP-09 & Injecție bibliotecă LD\_PRELOAD       & LSMFileWrite        & WARN     \\
\hline
REP-10 & Fork bomb                            & SyscallAnomaly      & WATCH    \\
\hline
REP-11 & Injecție prin ptrace                 & ProcessInjection    & FREEZE   \\
\hline
REP-12 & Brute-force SSH                      & NetworkAnomaly      & THROTTLE \\
\hline
\end{tabular}
\end{center}
\caption{Scenariile de atac din setul de bază și rezultatele așteptate}
\label{tab:basic_attacks}
\end{table}

\par REP-01 (rootkit Diamorphine) este implementat ca un modul de kernel inofensiv care nu produce efecte dăunătoare, dar este compilat și instalat prin \texttt{insmod}. Apelul \texttt{init\_module} sau \texttt{finit\_module} produce un scor de 95, declanșând imediat starea FREEZE. Acesta este scenariul cu cel mai scurt timp de detecție din suita de teste: evenimentul apare în ring buffer în milisecunde de la execuția \texttt{insmod}, înainte ca rootkit-ul să își poată ascunde prezența în kernel.

\par REP-11 (injecție ptrace) simulează tehnica clasică de deturnare a memoriei unui proces activ printr-o serie de apeluri \texttt{ptrace(PTRACE\_ATTACH, ...)}, \texttt{ptrace(PTRACE\_POKETEXT, ...)} și \texttt{ptrace(PTRACE\_DETACH, ...)}. Syscall-ul \texttt{ptrace} are ponderea de risc 75; amplificat cu 1.3 pentru root ($97.5$, clamped la 95), produce starea FREEZE.

\par REP-12 (brute-force SSH) generează conexiuni repetate la portul 22 printr-un program C care ciclează \texttt{socket()} + \texttt{connect()} de câteva sute de ori pe secundă. Stratul de burst detection detectează depășirea pragului de 5 apeluri \texttt{connect} în 3 secunde și emite un alert cu scor 75, producând starea THROTTLE.

\subsection{Scenariile avansate (REP-13 -- REP-21)}
\label{subsec:advanced_attacks}

\par Setul avansat cuprinde 9 scenarii cu tehnici mai sofisticate, prezentate în Tabelul~\ref{tab:advanced_attacks}.

\begin{table}[htbp]
\begin{center}
\begin{tabular}{|p{45pt}|p{110pt}|p{75pt}|p{55pt}|}
\hline
\textbf{ID} & \textbf{Scenariu} & \textbf{Tip alertă} & \textbf{Stare} \\
\hline\hline
REP-13 & Execuție fără fișier (memfd)           & MemoryInjection     & WARN     \\
\hline
REP-14 & Process hollowing (ptrace)             & ProcessInjection    & FREEZE   \\
\hline
REP-15 & Reverse shell                          & NetworkAnomaly      & THROTTLE \\
\hline
REP-16 & Recoltare credențiale (/etc/shadow)   & LSMFileWrite        & WARN     \\
\hline
REP-17 & Persistență finit\_module              & HiddenPID           & FREEZE   \\
\hline
REP-18 & Hook syscall (modul kernel)            & HiddenPID           & FREEZE   \\
\hline
REP-19 & Rootkit eBPF (stil TripleCross)        & ProcessInjection    & FREEZE   \\
\hline
REP-20 & Furt memorie proces                   & ProcessInjection    & FREEZE   \\
\hline
REP-21 & Tentativă VM escape                   & PrivilegeEscalation & ISOLATE  \\
\hline
\end{tabular}
\end{center}
\caption{Scenariile de atac din setul avansat și rezultatele așteptate}
\label{tab:advanced_attacks}
\end{table}

\par REP-13 (execuție fără fișier pe disc) simulează tehnica \textit{fileless malware}: codul este copiat într-o zonă de memorie anonimă creată cu \texttt{mmap(MAP\_ANONYMOUS)}, marcată executabilă cu \texttt{mprotect(PROT\_EXEC)}, și apelată direct. Combinația apelurilor \texttt{mmap} (pentru alocare) și \texttt{mprotect} (pentru acordarea drepturilor de execuție zonelor anonime) este detectată cu precizie de motorul logic, producând scoruri de 40, respectiv 45, și rezultând în starea WARN.

\par REP-19 (rootkit eBPF stil TripleCross) reproduce tehnica publicată de Pat Hogan (2022) \cite{Hogan2022}, în care un program eBPF malițios se atașează la apeluri de sistem pentru a modifica comportamentul proceselor. Implementarea din test utilizează un program eBPF inofensiv, dar apelează \texttt{ptrace} pentru a se atașa la un proces țintă, generând scorul de 95 și starea FREEZE.

\par REP-14 (process hollowing) necesită activarea \texttt{kernel.yama.ptrace\_scope=0} prin sysctl, deoarece Debian 12 restricționează implicit operațiunile ptrace la relații părinte-copil. Scenariul demonstrează că Sentinel detectează corect tehnica indiferent de configurația Yama, deoarece monitorizarea se face la nivelul syscall, nu la nivelul politicii Yama.


\section{Considerente de performanță}
\label{sec:performance}

\par Penalizarea de performanță a eBPF a fost evaluată comparând utilizarea CPU pe un workload de compilare kernel (\texttt{make -j4 bzImage}) cu și fără daemonul Sentinel activ, în mașina virtuală de test. Rezultatele medii pe 5 rulări sunt prezentate în Tabelul~\ref{tab:perf}.

\begin{table}[htbp]
\begin{center}
\begin{tabular}{|p{120pt}|p{80pt}|p{80pt}|}
\hline
\textbf{Configurație} & \textbf{Durată medie (s)} & \textbf{CPU user\%} \\
\hline\hline
Fără Sentinel                & 187.4 & 94.2 \\
\hline
Cu Sentinel (fără scorer ML) & 189.7 & 95.4 \\
\hline
Cu Sentinel (cu scorer ML)   & 192.1 & 96.1 \\
\hline
\end{tabular}
\end{center}
\caption{Comparație de performanță: compilare kernel cu și fără Sentinel activ}
\label{tab:perf}
\end{table}

\par Penalizarea de performanță introdusă de daemonul eBPF fără scorer ML este de $\approx 1.2\%$ din durata totală, sub pragul nefuncțional RN1 de 2\%. Scorerul ML adaugă $\approx 2.7$ secunde suplimentare, dar acesta operează asincron și nu blochează procesarea principală; impactul reflectă consumul de CPU al inferenței Python pe thread-uri separate.

\par Latența de carantinare (de la inserarea evenimentului în baza de date la aplicarea acțiunii) a fost măsurată prin câmpul \texttt{latency\_ms} din \texttt{quarantine\_log}. Pe cele 11 scenarii din setul de bază care produc stări $\geq$ THROTTLE, latența medie înregistrată a fost de $1.8$ ms pentru THROTTLE, $3.1$ ms pentru ISOLATE (include apelurile nftables) și $0.9$ ms pentru FREEZE (\texttt{SIGSTOP} este imediat, \texttt{cgroup.freeze} este asincron). Toate valorile se încadrează în cerința RN2 ($< 5$ ms).

\par În testul REP-10 (fork bomb, 100+ apeluri \texttt{clone}/s timp de 5 secunde), nu s-au înregistrat drop-uri în ring buffer-ul de 64 MiB. Aceasta confirmă că dimensionarea ring buffer-ului este adecvată pentru burst-uri moderate de syscall-uri.

\chapter{Concluzii}
\label{conclusions}

\par Lucrarea a prezentat proiectarea, implementarea și validarea sistemului Sentinel — un sistem de detectare și răspuns la intruziuni bazat pe eBPF pentru platforma Linux. Sentinel operează la nivel de kernel prin kprobes eBPF, calculează scoruri de risc printr-un motor hibrid cu trei straturi și aplică automat acțiuni graduale de containment prin cgroup v2 și nftables, cu o latență de răspuns măsurată sub 5 ms pe setul de scenarii de test.

\par Obiectivele formulate în Capitolul~\ref{intro} au fost atinse:

\begin{itemize}
    \item Programul eBPF (\texttt{tracer.bpf.c}) interceptează 21 de apeluri de sistem prin kprobes, cu o suprasarcină CPU de $\approx 1.2\%$ pe un workload de compilare intensivă — sub pragul RN1.
    \item Motorul de scoring hibrid combină ponderi statice ponderate, detectare a activității în rafale cu fereastră glisantă și un scorer Isolation Forest care poate fi antrenat cu date actualizate fără a întrerupe serviciul, producând o clasificare graduală în șase stări de severitate.
    \item Managerul de carantinare aplică restricții de CPU (cgroup \texttt{cpu.max}), izolare de rețea (nftables pe inodul cgroup-ului) și înghețare a proceselor (SIGSTOP + \texttt{cgroup.freeze}), cu auto-dezghețare temporizată.
    \item Suita de 21 de scenarii de atac reproductibile validează detecția corectă a categoriilor principale de amenințări, de la rootkit-uri de kernel până la execuție fără fișier și rootkit-uri eBPF.
\end{itemize}


\section*{Contribuții originale}
\addcontentsline{toc}{section}{Contribuții originale}

\par Contribuțiile originale față de instrumentele existente (Falco \cite{FalcoProject}, auditd \cite{LinuxAudit}) sunt:

\begin{enumerate}
    \item \textbf{Motorul de scoring cu trei straturi} cu combinare ponderată adaptivă: ponderea ML scade pe măsură ce riscul static crește, prevenind degradarea scorurilor pentru syscall-urile cu risc dovedit maxim.
    \item \textbf{Mecanismul de izolare de rețea bazat pe cgroup inode}: blocarea selectivă a traficului unui singur proces prin nftables + inodul cgroup-ului, fără a crea namespace-uri de rețea separate.
    \item \textbf{Cheia de burst cu start\_time}: prevenirea coliziunilor false în detectarea rafale pentru procese cu PID-uri reutilizate, prin incorporarea timpului de pornire al procesului în cheia de identificare.
    \item \textbf{Suita de 21 de atacuri reproductibile}: artefacte de testare publice care pot fi utilizate independent pentru validarea altor sisteme HIDS.
\end{enumerate}


\section*{Limitări}
\addcontentsline{toc}{section}{Limitări}

\par Sistemul prezintă o serie de limitări ce trebuie recunoscute explicit:

\begin{itemize}
    \item \textbf{Absența corelației între evenimente}: Sentinel evaluează fiecare eveniment independent. Un atac care distribuie activitatea suspectă între mai multe syscall-uri cu scoruri individuale scăzute poate evita detecția — o limitare fundamentală a abordărilor bazate pe evenimente individuale, față de modelele bazate pe secvențe \cite{Denning1987}.
    \item \textbf{Drift al distribuției ML}: Modelul Isolation Forest antrenat pe comportamentul normal inițial poate deveni neperformant dacă profilul de utilizare al sistemului se modifică semnificativ, necesitând antrenare periodică. Absența unui mecanism automat de detectare a drift-ului este o limitare practică.
    \item \textbf{Cerința de privilegii root}: Managerul de carantinare necesită privilegii root pentru manipularea cgroup-urilor și nftables; rularea fără privilegii reduce sistemul la detecție pură.
    \item \textbf{Acoperire limitată a izolării de rețea}: Blocarea prin nftables nu acoperă comunicarea prin loopback (\texttt{127.0.0.1}), permițând potențial lateral movement între procese de pe același host prin socket-uri locale.
\end{itemize}


\section*{Direcții de cercetare viitoare}
\addcontentsline{toc}{section}{Direcții de cercetare viitoare}

\par Limitările identificate trasează direcțiile principale de extindere:

\begin{itemize}
    \item \textbf{Corelarea secvențelor de syscall-uri}: Integrarea unui model de secvențe (Markov chain sau LSTM) pentru detectarea tiparelor de atac distribuite pe mai mulți pași — abordare propusă inițial de Denning \cite{Denning1987} și rafinată ulterior în literatura de ML pentru securitate \cite{Buczak2016}.
    \item \textbf{Suport uprobe}: Extinderea instrumentării la nivelul funcțiilor din spațiul utilizator (uprobe eBPF) pentru vizibilitate în biblioteci standard și interpretatoare — relevantă pentru detectarea atacurilor care operează exclusiv în spațiul utilizator.
    \item \textbf{Detecție automată a drift-ului}: Implementarea unui mecanism de detectare statistică a schimbărilor de distribuție (Kullback-Leibler divergence sau Jensen-Shannon) pentru declanșarea automată a antrenării cu date actualizate.
    \item \textbf{Extindere la sisteme distribuite}: Arhitectură multi-agent cu agregare centralizată a evenimentelor, permițând corelarea atacurilor ce traversează mai multe host-uri sau containere.
    \item \textbf{Suport CO-RE}: Refactorizarea programului eBPF pentru a utiliza BTF și relocări CO-RE, eliminând dependența de kernel headers la compilare și permițând distribuirea unui singur binar independent de versiunea kernel-ului.
\end{itemize}

\par Sentinelul demonstrează că eBPF oferă o fundație solidă pentru sistemele de securitate la nivel de kernel: vizibilitate completă, suprasarcină scăzută și integrare transparentă cu mecanismele existente ale sistemului de operare. Combinarea eBPF cu metode ML de detecție a anomaliilor și cu containment automat prin cgroup v2 produce un sistem capabil să detecteze și să limiteze impactul unui spectru larg de tehnici de atac documentate, fără a modifica kernel-ul sau a restricționa anticipat comportamentul proceselor.

%\addcontentsline{toc}{chapter}{Concluzii}
%\addcontentsline{toc}{chapter}{Conclusions}

\bibliography{references}

\end{document}
